\documentclass[a4paper,11pt]{article}
\usepackage[utf8]{inputenc}
\usepackage[english,francais]{babel}
\usepackage{amsmath}
\usepackage{bm}
\usepackage{amsfonts}
\usepackage{graphicx}
\usepackage{geometry}
\usepackage[utf8]{inputenc}
\usepackage{titlesec}
\usepackage{enumitem}
\usepackage{tikz}
\usepackage{calc}
\usepackage{caption}
\usepackage{easy-todo}
\usepackage{comment}
\usepackage{lipsum}
\usepackage{bbm}
\usepackage{fancyhdr}
\usepackage{tikz}
\usepackage{pdfpages}
\newcommand{\Var}{\mathbb{V}\text{ar}}
\newcommand{\Ex}{\mathbb{E}}
\newcommand{\Prob}{\mathbb{P}}

\newcommand{\ProbGP}{\mathcal{P}}
\newcommand{\Kspace}{\mathbb{K}}
\newcommand{\Uspace}{\mathbb{U}}
\newcommand{\Xspace}{\mathbb{X}}
\DeclareMathOperator{\Cov}{\mathrm{Cov}}
\DeclareMathOperator*{\argmin}{arg\,min \,}
\DeclareMathOperator*{\argmax}{arg\,max \,}
\DeclareMathOperator{\IMSE}{IMSE}
\usepackage{hyperref}
\usepackage{booktabs}

\pagestyle{fancy}
% \rhead{Victor Trappler}
% \lhead{}
% \graphicspath{{./Figures/}}
\begin{document}


\title{Notes}

\author{Victor Trappler \\[1cm]
  \begin{tabular}{lr}
    Directeurs de Thèse: & Arthur VIDARD (Inria) \\
                        & Élise ARNAUD (UGA)\\
                        & Laurent DEBREU (Inria)
  \end{tabular}
}

\maketitle
% \vspace{3cm}
% \includegraphics[scale=0.3]{/home/victor/logo_UGA}
% \hfill
% \includegraphics[scale=0.3]{/home/victor/ljk}
% \hfill
% \includegraphics[scale=0.3]{/home/victor/inria}
% \thispagestyle{empty} 
% \clearpage
\tableofcontents
\section{(Joint) Posterior formulation}
\subsection{Priors}
\begin{align*}
  K \sim \mathcal{U}(\mathbb{K}), \quad p(k) \\
  U \sim \mathcal{U}(\mathbb{U}), \quad p(u)
\end{align*}
\subsection{Likelihood model}
\begin{align*}
  p(y \mid k, u, \sigma^2) &= \frac{1}{\sqrt{2\pi}\sigma}\exp\left[-\frac{1}{2\sigma^2}SS(k,u)\right] \\
                         &= \frac{1}{\sqrt{2\pi}\sigma}\exp\left[-\frac{1}{2\sigma^2} \|\mathcal{M}(k,u) - y \|^2_{\Sigma}\right]
\end{align*}
Now to Bayes' theorem
\begin{align*}
  p(k,u \mid y,\sigma^2) = \frac{p(y \mid k, u, \sigma^2) p(k,u)}{\iint_{\mathbb{K}\times\mathbb{U}}p(y \mid k, u, \sigma^2) p(k,u) \, \mathrm{d}(k,u)}
\end{align*}
Let us assume an hyperprior for $\sigma^2$: $p(\sigma^2)$

\newpage
\section{GP, RR-based family of estimators}
\subsection{Random processes}
Let us assume that we have a map $f$ from a $p$ dimensional space to $\mathbb{R}$:
\begin{align}
  \begin{array}{rrcl}
    f: & \mathbb{X} \subset \mathbb{R}^p& \longrightarrow & \mathbb{R} \\
       & x & \longmapsto & f(x)
  \end{array}
\end{align}
This function is assumed to have been evaluated on a design of $n$ points, $\mathcal{X} \subset \mathbb{X}^n$. 
We wish to have a probabilistic modelling of this function
We introduce random processes as way to have a prior distribution on function
This uncertainty on $f$ is modelled as a random process:
\begin{equation}
  \begin{array}{rcl}
    Z: \mathbb{X} \times \Omega& \longrightarrow & \mathbb{R} \\
    (x,\omega) & \longmapsto & Z(x,\omega)
  \end{array}
\end{equation}
The $\omega$ variable will be omitted next.
\subsection{Linear Estimation}
\label{sec:linear_estimation}
A linear estimation $\hat{Z}$ of $f$ at an unobserved point $x\notin \mathcal{X}$ can be written as
\begin{equation}
  \label{eq:lin_est}
  \hat{Z}(x) =
  \begin{bmatrix}
    w_1 \dots w_n
    \end{bmatrix}
    \begin{bmatrix}
      f(x_1) \\ \vdots \\ f(x_n)
    \end{bmatrix} = \mathbf{W}^Tf(\mathcal{X}) = \sum_{i=1}^n w_i(x) f(x_i)
\end{equation}
Using those kriging weights $\mathbf{W}$, a few additional conditions must be added, in order to obtain the Best Linear Unbiased Estimator:
\begin{itemize}
\item Non-biased estimation: $\Ex[\hat{Z}(x) - Z(x)]=0$
\item Minimal variance: $\min~\Ex[(\hat{Z}(x) - Z(x))^2]$
\end{itemize}
Translating using Eq.\eqref{eq:lin_est}:
\begin{equation}
  \Ex[\hat{Z}(x) - Z(x)]=0 \iff m(\sum_{i=1}^n w_i(x)-1) = 0 \iff \sum_{i=1}^n w_i(x) = 1 \iff \mathbf{1}^T \mathbf{W} = 1
\end{equation}
For the minimum of variance, we introduce the augmented vector $\mathbf{Z}_n(x) = [Z(x_1),\dots Z(x_n), Z(x)]$, and
the variance can be expressed as:
\begin{align}
  \Ex[(\hat{Z}(x) - Z(x))^2] &= \Cov\left[[\mathbf{W}^T, -1] \cdot \mathbf{Z}_n(x) \right] \\
                             &= [\mathbf{W}^T, -1] \Cov\left[\mathbf{Z}_n(x) \right] [\mathbf{W}^T, -1]^T
\end{align}
In addition, we have
\begin{equation}
  \Cov\left[\mathbf{Z}_n(x) \right] =
  \begin{bmatrix}
    \Cov\left[
      \begin{bmatrix}
        Z(x_1) \dots Z(x_n)
      \end{bmatrix}^T\right]
    & \Cov\left[
      \begin{bmatrix}
        Z(x_1) \dots Z(x_n)
      \end{bmatrix}^T, Z(x) \right]
  \\
  \Cov\left[
    \begin{bmatrix}
      Z(x_1) \dots Z(x_n)
    \end{bmatrix}^T, Z(x) \right]^T & \Var\left[Z(x)\right]
  \end{bmatrix}
\end{equation}
Once expanded, the kriging weights solve then the following optimisation problem:
\begin{align}
  \min_{\mathbf{W}} ~&\mathbf{W}^T \Cov\left[Z(x_1) \dots Z(x_n)\right] \mathbf{W}\\ &-\Cov\left[
    \begin{bmatrix}
      Z(x_1) \dots Z(x_n)
    \end{bmatrix}^T, Z(x) \right]^T \mathbf{W}\\ &- \mathbf{W}^T\Cov\left[
    \begin{bmatrix}
      Z(x_1) \dots Z(x_n)
    \end{bmatrix}^T, Z(x) \right] \\ &+ \Var\left[Z(x)\right] \\
  \text{s.t.}& \mathbf{W}^T \mathbf{1} = \mathbf{1}
\end{align}
This leads to
\begin{align}
  \begin{bmatrix}
    \mathbf{W} \\ m
  \end{bmatrix}
  &=
  \begin{bmatrix}
    \Cov\left[Z(x_1) \dots Z(x_n)\right] & \mathbf{1} \\
  \mathbf{1}^T & 0
\end{bmatrix}^{-1}
                 \begin{bmatrix}
                  \Cov\left[
    \begin{bmatrix}
      Z(x_1) \dots Z(x_n)
    \end{bmatrix}^T, Z(x) \right]^T \\ 1 
\end{bmatrix}
  \\ &=
    \begin{bmatrix}
      C(x_1, x_1) & \cdots & C(x_1, x_n) & 1 \\
      C(x_2, x_1) & \cdots & C(x_2, x_n) & 1 \\
      \vdots & \ddots & \vdots & \vdots \\
      C(x_n, x_1) & \cdots & C(x_n, x_n)& 1 \\
      1 & \cdots & 1 & 0
    \end{bmatrix}^{-1}
                       \begin{bmatrix}
                         C(x_1, x) \\
                         C(x_2, x) \\
                         \vdots \\
                         C(x_n, x) \\
                         1
                       \end{bmatrix}
\end{align}
\subsection{GP of the penalized cost function $\Delta_{\alpha}$}
\subsubsection{GP processes}
Let $\Delta_{\alpha}(\mathbf{k},\mathbf{u}) = J(\mathbf{k},\mathbf{u}) - \alpha J^*(\mathbf{u})$. Furthermore, we assume that we constructed a GP on $J$ on the joint space $\Kspace \times \Uspace$, based on a design of $n$ points $\mathcal{X} = \left\{(\mathbf{k}^{(1)},\mathbf{u}^{(1)}),\dots,(\mathbf{k}^{(n)},\mathbf{u}^{(n)}) \right\}$, denoted as $(\mathbf{k},\mathbf{u})\mapsto Y(\mathbf{k},\mathbf{u})$.

As a GP, $Y$ is described by its mean function $m_{Y}$ and its covariance function $C(\cdot, \cdot)$, while $\sigma^2_Y(\mathbf{k},\mathbf{u}) = C((\mathbf{k},\mathbf{u}), (\mathbf{k},\mathbf{u}))$
\begin{equation}
  Y(\mathbf{k},\mathbf{u}) \sim \mathcal{N}\left(m_{Y}(\mathbf{k},\mathbf{u}), \sigma^2_Y(\mathbf{k},\mathbf{u}) \right)
\end{equation}
Let us consider now the conditional minimiser:
\begin{align}
  J^*(\mathbf{u}) = J(\mathbf{k}^*(\mathbf{u}),\mathbf{u}) = \min_{\mathbf{k}\in\Kspace} J(\mathbf{k},\mathbf{u})
\end{align}

Analogous to $J$ and $J^*$, we define $Y^*$ as
\begin{equation}
  Y^*(\mathbf{u}) \sim \mathcal{N}\left(m^*_Y(\mathbf{u}), \sigma^{2,*}_Y(\mathbf{u})\right)
\end{equation}
where
\begin{align}
  m^*_Y(\mathbf{u}) = \min_{\mathbf{k}\in\Kspace} m_Y(\mathbf{k},\mathbf{u})
\end{align}
The surrogate conditional minimiser is used in Ginsbourger profiles etc.
The $\alpha$-relaxed difference  $\Delta_{\alpha}$ modelled as a GP can then be written as

Considering the joint distribution of $Y(\mathbf{k},\mathbf{u})$ and $Y^*(\mathbf{u}) = Y(\mathbf{k}^*(\mathbf{u}), \mathbf{u})$, we have
\begin{equation}
  \begin{bmatrix}
    Y(\mathbf{k},\mathbf{u}) \\
    Y^*(\mathbf{u})
  \end{bmatrix}
  \sim \mathcal{N}\left(
    \begin{bmatrix}
      m_Y(\mathbf{k},\mathbf{u}) \\
      m_Y^*(\mathbf{u})
    \end{bmatrix}
    ;\,
    \begin{bmatrix}
      C\left((\mathbf{k},\mathbf{u}),(\mathbf{k},\mathbf{u})\right) & C\left((\mathbf{k},\mathbf{u}),(\mathbf{k}^*(\mathbf{u}),\mathbf{u})\right) \\
      C\left((\mathbf{k},\mathbf{u}),(\mathbf{k}^*(\mathbf{u}),\mathbf{u})\right) & C\left((\mathbf{k}^*(\mathbf{u}),\mathbf{u}),(\mathbf{k}^*(\mathbf{u}),\mathbf{u})\right)
    \end{bmatrix}
\right)
\end{equation}
By multiplying by the matrix $\left[1\quad -\alpha\right]$ yields
\begin{align}
  \Delta_{\alpha}(\mathbf{k},\mathbf{u}) &\sim \mathcal{N}\left(m_{\Delta}(\mathbf{k},\mathbf{u}); \sigma^2_{\Delta}(\mathbf{k},\mathbf{u})\right) \\
  m_{\Delta}(\mathbf{k},\mathbf{u}) &= m_Y(\mathbf{k},\mathbf{u}) - \alpha m_Y^*(\mathbf{u}) \\
  \sigma^2_{\Delta}(\mathbf{k},\mathbf{u}) &= \sigma_Y^2(\mathbf{k},\mathbf{u}) + \alpha^2 \sigma_{Y^*}^2(\mathbf{k},\mathbf{u}) - 2\alpha C\left((\mathbf{k},\mathbf{u}),(\mathbf{k}^*(\mathbf{u}),\mathbf{u})\right)
\end{align}
Assuming that $C((\mathbf{k},\mathbf{u}), (\mathbf{k}',\mathbf{u}')) = s \prod_{i\in\mathcal{I}_{\mathbf{k}}}\rho_{\theta_i}(\|k_i - k'_i\|) \prod_{j\in\mathcal{I}_{\mathbf{u}}} \rho_{\theta_j}(u_j - u'_j\|)$
\begin{align}
  C\left((\mathbf{k},\mathbf{u}),(\mathbf{k}^*(\mathbf{u}),\mathbf{u})\right) &= s \prod_{i\in\mathcal{I}_{\mathbf{k}}}\rho_{\theta_i}(\|k_i - k^*_i(\mathbf{u})\|)\prod_{j\in\mathcal{I}_{\mathbf{u}}} \rho_{\theta_j}(0) \\
  &=s \prod_{i\in\mathcal{I}_{\mathbf{k}}}\rho_{\theta_i}(\|k_i - k^*_i(\mathbf{u})\|)
\end{align}
\subsubsection{Approximation of the objective probability using GP}
We are going now to use a different notation for the probabilities, taken with respect to the GP: $\ProbGP$, to represent the uncertainty encompassed by the GP.

Defined somewhere else, we have
\begin{align}
  \Gamma_{\alpha}(\mathbf{k}) &= \Prob_{\mathbf{U}}\left[J(\mathbf{k},\mathbf{U}) \leq \alpha J^*(\mathbf{U})\right] \\
                              & =\Ex_{\mathbf{U}}\left[\mathbbm{1}_{J(\mathbf{k},\mathbf{U}) \leq \alpha J^*(\mathbf{U})}\right]
\end{align}
This classification problem can be approached with a plug-in approach, or a probablistic one:
\begin{align}
  \mathbbm{1}_{J(\mathbf{k},\mathbf{u}) \leq \alpha J^*(\mathbf{u})} &\approx   \mathbbm{1}_{m_Y(\mathbf{k},\mathbf{u}) \leq \alpha m_Y^*(\mathbf{u})} \\
  \mathbbm{1}_{J(\mathbf{k},\mathbf{u}) \leq \alpha J^*(\mathbf{u})} &\approx   \ProbGP\left[ \Delta_{\alpha}(\mathbf{k},\mathbf{u}) \leq 0 \right] = \pi(\mathbf{k},\mathbf{u})
\end{align}
Using the GPs, for a given $\mathbf{k}$, $\alpha$ and $\mathbf{u}$, the probability for our meta model to verify the inequality is given by
Based on those two approximation, the approximated probability $\Gamma$ is
\begin{align}
  \hat{\Gamma}_{\alpha, n}(\mathbf{k}) &= \Prob_U\left[m_Y(\mathbf{k},\mathbf{u}) \leq \alpha m_Y^*(\mathbf{u}) \right] \tag{plug-in} \\
  \hat{\Gamma}_{\alpha, n}(\mathbf{k}) &= \Ex_U\left[ \ProbGP\left[ \Delta_{\alpha}(\mathbf{k},\mathbf{u}) \leq 0\right]\right] \tag{Probabilistic approx} \\
\end{align}

The probability of coverage for the set $\{Y - \alpha Y^*\}$ is $\pi_{\alpha}$, and can be computed using the CDF of the standard normal distribution $\Phi$:
\begin{equation}
  \pi_{\alpha}(\mathbf{k},\mathbf{u}) = \Phi\left(-\frac{m_{\Delta_\alpha}(\mathbf{k},\mathbf{u})}{\sigma_{\Delta_\alpha}(\mathbf{k},\mathbf{u})}\right)
\end{equation}
Finally, averaging over $\mathbf{u}$ yields
\begin{equation}
  \hat{\Gamma}(\mathbf{k}) = \int_{\Uspace}\pi_{\alpha}(\mathbf{k},\mathbf{u})p(\mathbf{u}) \,\mathrm{d}\mathbf{u}
\end{equation}
\subsection{Sources, quantification of uncertainties, and SUR strategy ?}

Formally, for a given point $(\mathbf{k},\mathbf{u})$, the event ``the point is $\alpha$-acceptable'' has probability $\pi(\mathbf{k},\mathbf{u})$ and variance $\pi(\mathbf{k},\mathbf{u}) (1-\pi(\mathbf{k},\mathbf{u}))$. Obviously, the points with the highest uncertainty have the highest variance, so have a coverage probability $\pi$ around 0.5.

\subsubsection{Random sets}
Let us start by introducing diverse tools based around Vorob'ev expectation of closed sets (ref thèse Reda), \cite{heinrich_level_2012}.


Let us consider $A$, a random closed set, such that its realizations are subsets of $\Xspace$, and $p$ is its coverage probability, that is
\begin{equation}
  p(x) = \Prob\left[x\in A\right], x\in\Xspace
\end{equation}
For $\eta \in [0, 1]$, we define the $\eta$-level set of $p$,
\begin{equation}
  Q_{\eta} = \{x\in\Xspace \mid p(x) \geq \eta \}
\end{equation}
It may seem trivial, but let us still note that those sets are decreasing:
\begin{equation}
  0\leq \eta \leq \xi \leq 1 \implies Q_{\xi} \subseteq Q_{\eta}
\end{equation}

Using those level sets for the level $\eta=0.05$ for instance:
\begin{equation}
  Q_{1-\frac{\eta}{2}} \subset Q_{\frac{\eta}{2}}
\end{equation}
Recalling the objective, it gives upper bounds and lower bounds of the confidence interval of level $\eta$ on the probability for each $\mathbf{k}$:
\begin{align}
  \hat{\Gamma_{\alpha}}^{U}(\mathbf{k}) &= \Prob_\mathbf{U}\left[x=(\mathbf{k},\mathbf{u}) \in Q_{1-\frac{\eta}{2}}\right] \\
  \hat{\Gamma_{\alpha}}^{L}(\mathbf{k}) &= \Prob_\mathbf{U}\left[x=(\mathbf{k},\mathbf{u}) \in Q_{\frac{\eta}{2}}\right]
\end{align}
In \cite{dubourg_reliability-based_2011} is introduced the Margin of uncertainty, defined as the following set difference
\begin{equation}
  \mathbb{M}_{\eta} = Q_{\frac{\eta}{2}} \setminus Q_{1-\frac{\eta}{2}}
\end{equation}
Considering the 

Let $\mu$ be a Borel $\sigma$-finite measure on $\Xspace$. We define Vorob'ev expectation, as the $\eta^*$-level set of $A$ verifying
\begin{equation}
  \forall \beta < \eta^* \quad \mu(Q_{\beta}) \leq \Ex[\mu(A)] \leq \mu(Q_{\eta^*})
\end{equation}
that is the level set of $p$, that has the volume of the mean of the volume of the random set $A$.

\subsubsection{SUR Strategies}
The main idea behind Stepwise Uncertainty Reduction is to define a criterion, say $\kappa_n$, that encapsulates the epistemic uncertainty, and to minimize this criterion, in order to select the next point:
\begin{equation}
  x^{n+1} = \argmax_{x\in\Xspace} \kappa_n(x)
\end{equation}
where $\kappa_n$ depends on $Y\mid \mathcal{X}_n$
This approach is suitable for step by step evaluations.

\subsubsection{Integrated Mean square criterion}
\cite{sacks_designs_1989}
Let us consider that we have a kriging model over $\Xspace$ based on a experimental design $\mathcal{X}$, that is denoted $Y \mid \mathcal{X}$

We define the Integrated Mean Square Error (IMSE) as
\begin{equation}
  \IMSE(Y \mid \mathcal{X}) = \int_{\Xspace} \sigma_{Y\mid\mathcal{X}}^2(x)\,\mathrm{d}x
\end{equation}
where
\begin{equation}
  Y\mid \mathcal{X} \sim \mathcal{N}(m_{Y\mid\mathcal{X}}(x),\sigma^2_{Y\mid\mathcal{X}}(x))
\end{equation}


\begin{equation}
  x^{n+1} = \argmin_{x\in \Xspace}\Ex_{y\sim Y(x)}\left[\IMSE\left(Y \mid \mathcal{X}\cup \left\{(x, y)\right\}\right) \right]
\end{equation}
So we choose th point minimizing the expected integrated mean square error.

\subsubsection{Weighted IMSE}
To include a more precise objective than the enrichment of the design, one can add a weight function to the integral, giving the $W-\IMSE$:
\begin{equation}
  \label{eq:w-imse}
  w-\IMSE(Y\mid \mathcal{X}) = \int_{\Xspace} \sigma_{Y\mid\mathcal{X}}^2(x)w(x)\,\mathrm{d}x
\end{equation}

In order to increase the accuracy of the surrogate model around some region of interest, the $w-\IMSE$ can be transformed into
\begin{equation}
  w-\IMSE(Y\mid \mathcal{X}) = \int_{\Xspace} \sigma_{Y\mid\mathcal{X}}^2(x)\ProbGP\left[x \in \mathbb{M}_{\eta}\right]\,\mathrm{d}x
\end{equation}
where $\mathcal{M}_{\eta}$ is the $\eta$-margin of uncertainty.

\subsubsection{UB-LB for $(p, \alpha_p, \mathbf{k}_p)$}
Let us assume that we have set a probability $p\in [0,1]$. Let us recall that the triplet $(p, \alpha_p, \mathbf{k}_p)$ verifies
\begin{align}
  \max_{\mathbf{k}} \Gamma_{\alpha_p}(\mathbf{k}) = \Gamma_{\alpha_p}(\mathbf{k}_p) = \Prob_{\mathbf{U}}\left[J(\mathbf{k}_p,\mathbf{U}) \leq \alpha_p J^*(\mathbf{U})\mid \mathbf{U} = \mathbf{u}\right] = p
\end{align}
Let us say that $\bar{\Gamma}$ is the $\eta$-upper-bound, while $\underline{\Gamma}$ is the $\eta$-lower bounds, so
\begin{equation}
  \ProbGP\left[\underline{\Gamma}(\mathbf{k}) \leq \Gamma_n(\mathbf{k}) \leq \bar{\Gamma}(\mathbf{k})\right] = \eta
\end{equation}
\begin{itemize}
\item If $\underline{\Gamma}(\mathbf{k})>p$, we are too permissive, so we should decrease $\alpha$
  \begin{itemize}
  \item by how much ?
  \end{itemize}
\item If $\bar{\Gamma}(\mathbf{k})<p$, we are too conservative, so we should increase $\alpha$
  \begin{itemize}
  \item by how much again ?
  \end{itemize}
 \item If $\underline{\Gamma}(\mathbf{k})<p<\bar{\Gamma}(\mathbf{k})$, reduce uncertainty on $\mathbf{k}_p$
\end{itemize}
Changing the value of $\alpha$ does not require any further evaluation of the objective function, so can be increased until $\max \hat{\Gamma} = p$ ? by dichotomy for instance. This $\hat{\mathbf{k}}_p$ is then the candidate.

Criterion: stepwise reduction of the variance of the estimation of $\hat{\Gamma}(\hat{\mathbf{k}}_p) = \max_{\mathbf{k}}\hat{\Gamma}(\hat{\mathbf{k}})$

For a fixed $p\in (0, 1]$, and an initial design $\mathcal{X}$. Set an initial value for $\alpha \geq 1$. 
\begin{itemize}
\item Define $\Delta_{\alpha}$, using $Y \mid \mathcal{X}$
\item Update $\alpha$ such that $\max \hat{\Gamma}_{\alpha,n} = p$
\item Compute measure of uncertainty that we want to reduce:
  \begin{itemize}
  \item $\bar{\Gamma}_{\alpha,n}(\mathbf{k}) - \underline{\Gamma}_{\alpha,n}(\mathbf{k})$
  \item $\pi_{\alpha}(\mathbf{k},\mathbf{u})(1-\pi_{\alpha}(\mathbf{k},\mathbf{u}))$
  \end{itemize}
\end{itemize}



\bibliographystyle{alpha}
\bibliography{/home/victor/acadwriting/bibzotero}
\end{document}
%%% Local Variables:
%%% mode: latex
%%% TeX-master: t
%%% End:
